\chapter{Evaluation and Results}
\label{ch:results}

This chapter presents our experimental setup and evaluation results. We compare the performance of RubricNet against baseline classifiers across three descriptor generations (V1, V2, and V3) and present results from our coarse classification and auxiliary model tuning experiments.

\section{Experimental Design}
To evaluate the models robustly, we implement a fixed, stratified 5-fold cross-validation split based on the binned GuitarBurst difficulty labels. This split (defined in \texttt{guitar\_splits.json}) is held invariant across all models and runs to prevent data leakage and ensure comparable evaluation.

We report six primary metrics:
\begin{enumerate}
    \item \textbf{Accuracy (Exact)}: The ratio of correctly predicted difficulty bins.
    \item \textbf{Balanced Accuracy}: The average of recall obtained on each class, correcting for class size imbalances.
    \item \textbf{Accuracy $\pm$ 1}: The fraction of predictions within $\pm 1$ bin of the true class (reported for informational purposes).
    \item \textbf{Mean Absolute Error (MAE)}: The average absolute difference between predicted and true bin indices.
    \item \textbf{Mean Squared Error (MSE)}: The average squared difference, penalizing larger prediction errors.
    \item \textbf{Kendall's $\tau$}: A rank correlation coefficient measuring ordinal alignment.
\end{enumerate}

All scores are reported as the mean $\pm$ standard deviation across the 5 cross-validation folds.

\section{8-Class Difficulty Estimation Results}
Table~\ref{tab:8class_results} presents the comparative evaluation of the models on the 8-class classification task.

\begin{table}[htbp]
\centering
\caption{8-Class Difficulty Estimation Comparison (Mean $\pm$ Std across 5 Folds)}
\label{tab:8class_results}
\resizebox{\textwidth}{!}{
\begin{tabular}{lcccccc}
\toprule
\textbf{Model} & \textbf{Accuracy} & \textbf{Balanced Acc} & \textbf{Acc $\pm$ 1} & \textbf{MAE} & \textbf{MSE} & \textbf{Kendall $\tau$} \\
\midrule
Ordinal Regression V1 & 0.214 $\pm$ 0.062 & 0.197 $\pm$ 0.051 & N/A & 1.447 $\pm$ 0.260 & 3.431 $\pm$ 0.907 & N/A \\
Decision Tree V1 & 0.282 $\pm$ 0.030 & 0.255 $\pm$ 0.033 & N/A & 1.352 $\pm$ 0.086 & 3.473 $\pm$ 0.338 & N/A \\
Random Forest V1 & 0.316 $\pm$ 0.041 & 0.278 $\pm$ 0.026 & N/A & 1.222 $\pm$ 0.081 & 2.951 $\pm$ 0.223 & N/A \\
RubricNet V1 & 0.243 $\pm$ 0.042 & 0.231 $\pm$ 0.031 & N/A & 1.258 $\pm$ 0.092 & 2.820 $\pm$ 0.380 & N/A \\
\midrule
Ordinal Regression V2 & 0.218 $\pm$ 0.056 & 0.214 $\pm$ 0.056 & N/A & 1.316 $\pm$ 0.339 & 2.906 $\pm$ 1.172 & N/A \\
Decision Tree V2 & 0.286 $\pm$ 0.036 & 0.248 $\pm$ 0.031 & N/A & 1.295 $\pm$ 0.118 & 3.128 $\pm$ 0.550 & N/A \\
Random Forest V2 & \textbf{0.320 $\pm$ 0.029} & 0.275 $\pm$ 0.025 & N/A & 1.116 $\pm$ 0.047 & 2.487 $\pm$ 0.202 & N/A \\
RubricNet V2 (Ours) & 0.300 $\pm$ 0.032 & \textbf{0.283 $\pm$ 0.033} & 0.730 $\pm$ 0.052 & \textbf{1.115 $\pm$ 0.102} & \textbf{2.372 $\pm$ 0.445} & \textbf{0.623 $\pm$ 0.038} \\
\midrule
Ordinal Regression V3 & 0.309 $\pm$ 0.033 & 0.288 $\pm$ 0.033 & N/A & 1.109 $\pm$ 0.143 & 2.467 $\pm$ 0.762 & N/A \\
Decision Tree V3 & 0.279 $\pm$ 0.036 & 0.242 $\pm$ 0.034 & N/A & 1.286 $\pm$ 0.059 & 3.130 $\pm$ 0.491 & N/A \\
Random Forest V3 & \textbf{0.331 $\pm$ 0.022} & 0.289 $\pm$ 0.019 & N/A & 1.122 $\pm$ 0.071 & 2.527 $\pm$ 0.312 & N/A \\
RubricNet V3 (Ours) & 0.310 $\pm$ 0.027 & \textbf{0.291 $\pm$ 0.028} & 0.733 $\pm$ 0.032 & \textbf{1.069 $\pm$ 0.063} & \textbf{2.147 $\pm$ 0.255} & \textbf{0.634 $\pm$ 0.026} \\
\bottomrule
\end{tabular}
}
\end{table}

\subsection{Performance Evolution Analysis}
The results show a clear upward trend across descriptor generations. 
In V1, RubricNet underperformed relative to the Random Forest baseline, lagging by 7.3\% in exact accuracy. This was due to corrupted fret values in the PDF parser and the presence of dead features.

In V2, after fixing the parser and introducing fret height and stretch distribution features, RubricNet's performance increased significantly, achieving 28.3\% balanced accuracy (outperforming Random Forest V2 on balanced accuracy, MAE, MSE, and Kendall's $\tau$).

In V3, with the inclusion of windowed bottleneck and rhythm-aware context features, RubricNet V3 achieved its highest performance: \textbf{31.0\% exact accuracy} and \textbf{29.1\% balanced accuracy}, outperforming Random Forest V3 on balanced accuracy, MAE, MSE, and Kendall's $\tau$ while remaining a fully interpretable additive model.

\section{Coarse 3-Class Evaluation}
To align with other publications that classify difficulty into three high-level ranges (Easy, Medium, Hard), we mapped our 8-class predictions at evaluation time into three classes. The mapping groupings are: Classes 0–2 $\to$ Easy, Classes 3–5 $\to$ Medium, and Classes 6–7 $\to$ Hard. The results are summarized in Table~\ref{tab:3class_results}.

\begin{table}[h]
\centering
\caption{Coarse 3-Class Evaluation (Easy / Medium / Hard)}
\label{tab:3class_results}
\begin{tabular}{lcc}
\toprule
\textbf{Model} & \textbf{Accuracy} & \textbf{Balanced Accuracy} \\
\midrule
RubricNet V2 (Coarse) & 0.670 $\pm$ 0.039 & 0.621 $\pm$ 0.046 \\
RubricNet V3 (Coarse) & 0.676 $\pm$ 0.033 & 0.614 $\pm$ 0.052 \\
\bottomrule
\end{tabular}
\end{table}

Under the coarse evaluation scheme, RubricNet V3 achieves an exact accuracy of \textbf{67.6\%}, proving its practical utility for general repertoire categorization.

\section{Auxiliary Experimental Findings}

\subsection{Feature Pruning}
We conducted an experiment where we dropped the 6 weakest features identified in the Spearman correlation audit (e.g., \texttt{tempo\_bpm}, \texttt{arpeggio\_density}, \texttt{fret\_change\_rate}, \texttt{avg\_string\_jump}, \texttt{chord\_ratio}, \texttt{avg\_polyphony}). Interestingly, training RubricNet on this pruned set *reduced* accuracy (from 31.0\% to 30.2\%) and increased variance. This suggests that RubricNet's Optuna-tuned regularizers (dropout and weight decay) are already effective at neutralizing uninformative inputs, while retaining them provides minor stabilization.

\subsection{Ordinal Label Smoothing (OLS)}
Sweeping the smoothing temperature $T \in \{0.15, 0.3, 0.6\}$ on RubricNet V3 showed that a low temperature of \textbf{$T=0.15$} yields a small, consistent improvement (accuracy increased to 31.2\% and balanced accuracy to 29.4\%). High smoothing temperatures ($T \ge 0.3$) degraded performance, as over-smoothing erodes the ordinal signal.

\subsection{Multi-Task Auxiliary Head}
We trained a model with a joint, coarse-class auxiliary head to predict both the 8-class and 3-class targets simultaneously. This model performed worse across all folds (accuracy 29.3\%–30.7\%). Because RubricNet restricts its shared trunk to a single aggregated scalar, the model lacks the capacity to optimize for two conflicting sets of class boundaries simultaneously.
